\section{Step 48 Keep file identity separate from the replay cursor}
Fill \texttt{seek\_state}. Retain the directory pointer in \texttt{\$s4} and scratch pointer in \texttt{\$s7}. Set the search/replay index \texttt{\$s5} to the requested target \texttt{\$s1}. The supplied frame has already retained the file base, target, destination and context in \texttt{\$s0} through \texttt{\$s3}. The target remains fixed while the search index moves backward. Later nested record calls may replace temporary registers; keep values that must survive in saved registers or private storage.

Use \texttt{s48-1} and stop before \texttt{seek\_state\_ready}.

\paragraph{Observe} At \texttt{seek\_state\_ready}, \texttt{\$s5} is six and \texttt{\$s7} names scratch. Neither value is the currently displayed frame index.

\paragraph{Check} Which value must remain fixed during the backward search?
\begin{enumerate}[label=\Alph*.]
\item The current directory index.
\item The current record address.
\item The requested target.
\end{enumerate}
